% appP_graph_b.tex -- generated by writing/tools/blueprint_gen.py 1.0 (level 0) on 2026-10-14 from sources 20021204b9; do not edit
\providecommand{\bpNum}[1]{\hyperref[#1]{\ref*{#1}}}%
\providecommand{\bpNumA}[1]{\hyperref[#1]{\textcolor{orange!65!black}{\ref*{#1}}}}%
\providecommand{\bpPage}[1]{\pageref*{#1}}%
\providecommand{\bpLab}[1]{{\ttfamily #1}}%
\providecommand{\bpBrk}{\allowbreak}%
\providecommand{\bpT}[2]{\facttagbox[#1]{\scriptsize #2}}%
\providecommand{\bpU}[1]{{\color{gray!70!black}#1}}%
\begin{figure}[p]
\centering
\rotatebox{90}{\begin{tikzpicture}[
  bp node/.style={draw,line width=0.45pt,rounded corners=1pt,font=\tiny\ttfamily,align=center,
                  inner xsep=1.5pt,inner ysep=1.1pt,minimum height=4mm},
  bp R/.style={bp node,fill=OliveGreen!15,draw=OliveGreen!60!black},
  bp A/.style={bp node,fill=RoyalBlue!15,draw=RoyalBlue!60!black},
  bp P/.style={bp node,fill=BurntOrange!15,draw=BurntOrange!70!black},
  bp C/.style={bp node,fill=Plum!15,draw=Plum!70!black},
  bp V/.style={bp node,fill=BrickRed!15,draw=BrickRed!70!black,double,double distance=0.6pt},
  bp cond/.style={dash pattern=on 1.6pt off 1pt},
  bp math/.style={bp node,fill=white,draw=black!60},
  bp def/.style={bp node,fill=white,draw=black!45,line width=0.3pt,rounded corners=0pt},
  bp conj/.style={bp node,ellipse,inner xsep=0pt,inner ysep=0pt,fill=Plum!15,draw=Plum!70!black},
  bp asm/.style={bp node,fill=orange!8,draw=orange!65!black,dash pattern=on 2pt off 1.2pt},
  bp stub/.style={bp node,font=\tiny,fill=gray!12,draw=gray!55,text=black!70},
  bp fold/.style={bp node,font=\tiny,fill=white,draw=gray!55,dash pattern=on 1pt off 1pt,text=black!70},
  bp unit/.style={font=\scriptsize,text=black!70,inner sep=1pt},
  bp sep/.style={draw=gray!40,line width=0.4pt,dash pattern=on 1pt off 1.5pt},
  bp key/.style={font=\scriptsize,anchor=west,inner sep=0pt},
  bp use/.style={draw=gray!60,line width=0.5pt,-{Stealth[length=2.6pt,width=2.2pt]}},
  bp stubuse/.style={draw=gray!45,line width=0.45pt,dash pattern=on 1.2pt off 1pt,-{Stealth[length=2.6pt,width=2.2pt]}},
  bp asmuse/.style={draw=orange!65!black,line width=0.5pt,dash pattern=on 2pt off 1.2pt,-{Stealth[length=2.6pt,width=2.2pt]}},
  bp neck/.style={draw=red!60!black,line width=1.1pt,-{Stealth[length=3.6pt,width=3.2pt]}},
  bp long/.style={opacity=0.5}]
  \node[bp unit] at (1.792,5.262) {第~2~章};
  \node[bp unit] at (7.010,5.262) {第~3~章};
  \node[bp unit] at (11.538,5.262) {第~4~章};
  \node[bp unit] at (16.712,5.262) {第~5~章};
  \draw[bp sep] (4.775,5.122) -- (4.775,-4.992);
  \draw[bp sep] (9.072,5.122) -- (9.072,-4.992);
  \draw[bp sep] (14.176,5.122) -- (14.176,-4.992);
  \draw[bp sep] (19.101,5.122) -- (19.101,-4.992);
  \node[bp R] (n1) at (0.000,1.693) {\detokenize{qcd-}\\\detokenize{covariance}};
  \node[bp R] (n2) at (1.648,2.241) {\detokenize{qcd-gauge-}\\\detokenize{invariance}};
  \node[bp R] (n3) at (0.000,2.766) {\detokenize{qcd-}\\\detokenize{superficial-}\\\detokenize{degree}};
  \node[bp R] (n4) at (1.648,1.145) {\detokenize{prelim-brst-}\\\detokenize{nilpotent}};
  \node[bp R] (n5) at (3.583,-0.154) {\detokenize{qcd-brst-}\\\detokenize{cohomology}};
  \node[bp R] (n6) at (3.583,1.875) {\detokenize{qcd-st-}\\\detokenize{regularized}};
  \node[bp R] (n7) at (3.583,2.766) {\detokenize{qcd-}\\\detokenize{renormalizability}};
  \node[bp R] (n8) at (3.583,4.621) {\detokenize{qcd-st-}\\\detokenize{relations}};
  \node[bp R] (n9) at (3.583,0.860) {\detokenize{prelim-}\\\detokenize{slavnov-}\\\detokenize{taylor}};
  \node[bp R] (n10) at (3.583,-1.046) {\detokenize{qcd-xi-}\\\detokenize{independence}};
  \node[bp R] (n11) at (3.583,-2.060) {\detokenize{prelim-}\\\detokenize{composite-}\\\detokenize{renorm}};
  \node[bp R] (n12) at (3.583,-2.954) {\detokenize{prelim-beta}};
  \node[bp R] (n13) at (3.583,-3.604) {\detokenize{qcd-rge}};
  \node[bp R] (n14) at (3.583,-4.498) {\detokenize{qcd-}\\\detokenize{asymptotic-}\\\detokenize{freedom}};
  \node[bp R] (n15) at (5.679,2.383) {\detokenize{prelim-boost}};
  \node[bp R] (n16) at (5.679,0.475) {\detokenize{prelim-}\\\detokenize{wilson-rules}};
  \node[bp math] (n17) at (5.679,1.489) {\detokenize{prelim-}\\\detokenize{eikonal-}\\\detokenize{identity}};
  \node[bp math] (n18) at (5.679,-0.417) {\detokenize{prelim-}\\\detokenize{scaleless}};
  \node[bp R] (n19) at (7.096,0.224) {\detokenize{prelim-}\\\detokenize{cutting-}\\\detokenize{rules}};
  \node[bp R] (n20) at (8.341,0.038) {\detokenize{prelim-}\\\detokenize{unitarity}};
  \node[bp math] (n21) at (9.919,3.526) {\detokenize{ir-envelope}};
  \node[bp R] (n22) at (11.451,3.526) {\detokenize{ir-landau-}\\\detokenize{alpha}};
  \node[bp R] (n23) at (9.919,1.294) {\detokenize{ir-landau-}\\\detokenize{momentum}};
  \node[bp R] (n24) at (11.451,2.635) {\detokenize{ir-coleman-}\\\detokenize{norton}};
  \node[bp A] (n25) at (13.157,2.622) {\detokenize{ir-}\\\detokenize{pinch-}\\\detokenize{classification}};
  \node[bp R] (n26) at (9.919,2.186) {\detokenize{ir-power-}\\\detokenize{counting}};
  \node[bp A] (n27) at (11.451,1.743) {\detokenize{ir-super-}\\\detokenize{leading}};
  \node[bp A] (n28) at (13.157,1.180) {\detokenize{ir-leading-}\\\detokenize{regions}};
  \node[bp R] (n29) at (11.451,0.852) {\detokenize{ir-glauber-}\\\detokenize{trapping}};
  \node[bp R] (n30) at (9.919,-1.208) {\detokenize{ir-bn-qed}};
  \node[bp A] (n31) at (11.451,-0.810) {\detokenize{ir-kln}};
  \node[bp V] (n32) at (11.451,-1.606) {\detokenize{ir-bn-}\\\detokenize{failure}};
  \node[bp A] (n33) at (13.157,0.029) {\detokenize{ir-safe-}\\\detokenize{finite}};
  \node[bp R] (n34) at (9.919,-0.260) {\detokenize{ir-regions-}\\\detokenize{euclid}};
  \node[bp P] (n35) at (11.451,-0.039) {\detokenize{ir-regions-}\\\detokenize{minkowski}};
  \node[bp R] (n36) at (15.080,0.074) {\detokenize{sudakov-}\\\detokenize{oneloop}};
  \node[bp A] (n37) at (16.785,0.629) {\detokenize{sudakov-}\\\detokenize{factorization}};
  \node[bp A] (n38) at (15.080,1.012) {\detokenize{sudakov-webs}};
  \node[bp A] (n39) at (18.344,0.799) {\detokenize{sudakov-}\\\detokenize{ir-poles}};
  \node[bp V] (n40) at (18.344,-0.329) {\detokenize{sudakov-}\\\detokenize{dipole-}\\\detokenize{violation}};
  \node[bp stub] (n41) at (19.648,1.413) {附录 B};
  \node[bp stub] (n42) at (19.648,0.763) {附录 D};
  \node[bp stub] (n43) at (19.648,2.129) {附录 E};
  \node[bp stub] (n44) at (19.648,0.029) {附录 J};
  \begin{scope}[on background layer]
  \draw[bp use] (n1) -- (n2);
  \draw[bp use] (n1) -- (n4);
  \draw[bp use] (n2) -- (n7);
  \draw[bp use] (n3) -- (n7);
  \draw[bp use] (n4) -- (n5);
  \draw[bp use] (n4) -- (n6);
  \draw[bp use] (n5) to[bend left=35] (n10);
  \draw[bp use,bp long] (n5) -- (n20);
  \draw[bp use] (n6) to[bend left=35] (n7);
  \draw[bp use] (n7) to[bend left=35] (n8);
  \draw[bp use] (n7) to[bend left=35] (n9);
  \draw[bp use] (n7) to[bend left=35] (n11);
  \draw[bp use] (n8) to[bend left=35] (n12);
  \draw[bp use] (n9) to[bend left=35] (n10);
  \draw[bp use,bp long] (n9) -- (n20);
  \draw[bp use,bp long] (n9) -- (n27);
  \draw[bp use] (n10) to[bend left=35] (n13);
  \draw[bp use] (n12) to[bend left=35] (n13);
  \draw[bp use] (n13) to[bend left=35] (n14);
  \draw[bp stubuse,bp long] (n41) -- (n16);
  \draw[bp stubuse,bp long] (n41) -- (n17);
  \draw[bp stubuse,bp long] (n41) -- (n38);
  \draw[bp stubuse,bp long] (n42) -- (n23);
  \draw[bp use,bp long] (n15) -- (n26);
  \draw[bp use] (n16) -- (n19);
  \draw[bp use,bp long] (n16) -- (n37);
  \draw[bp use,bp long] (n16) -- (n38);
  \draw[bp use,bp long] (n17) -- (n37);
  \draw[bp use,bp long] (n17) -- (n38);
  \draw[bp use,bp long] (n18) -- (n34);
  \draw[bp use,bp long] (n18) -- (n36);
  \draw[bp use] (n19) -- (n20);
  \draw[bp use] (n20) -- (n31);
  \draw[bp use,bp long] (n20) -- (n33);
  \draw[bp stubuse,bp long] (n43) -- (n26);
  \draw[bp stubuse,bp long] (n44) -- (n20);
  \draw[bp use] (n30) -- (n31);
  \draw[bp use] (n30) -- (n32);
  \draw[bp use] (n34) -- (n35);
  \draw[bp use,bp long] (n35) -- (n37);
  \draw[bp use] (n21) -- (n22);
  \draw[bp use] (n23) -- (n29);
  \draw[bp use] (n23) -- (n24);
  \draw[bp use] (n24) -- (n25);
  \draw[bp use] (n25) to[bend left=35] (n28);
  \draw[bp use] (n26) -- (n27);
  \draw[bp use] (n27) -- (n28);
  \draw[bp use] (n28) to[bend left=35] (n33);
  \draw[bp use] (n28) -- (n37);
  \draw[bp use] (n36) -- (n37);
  \draw[bp use] (n37) -- (n39);
  \draw[bp use] (n38) -- (n39);
  \draw[bp use] (n39) to[bend left=35] (n40);
  \end{scope}
  \begin{scope}[shift={(-0.754,-5.492)}]
    \node[bp R,anchor=west,minimum width=4.5mm] (k1) at (0,0) {R};
    \node[bp R,bp cond,anchor=west,minimum width=4.5mm,right=0.8mm of k1] (k2) {R};
    \node[bp A,anchor=west,minimum width=4.5mm,right=0.8mm of k2] (k3) {A};
    \node[bp P,anchor=west,minimum width=4.5mm,right=0.8mm of k3] (k4) {P};
    \node[bp C,anchor=west,minimum width=4.5mm,right=0.8mm of k4] (k5) {C};
    \node[bp V,anchor=west,minimum width=4.5mm,right=0.8mm of k5] (k6) {V};
    \node[bp key,right=1.2mm of k6] (t1) {严格性标签（虚线框为 $R_{\rm cond}$）};
    \node[bp conj,anchor=west,minimum width=7mm,minimum height=3.4mm,right=2.5mm of t1] (k7) {};
    \node[bp key,right=1mm of k7] (t2) {猜想};
    \node[bp asm,anchor=west,minimum width=4.5mm,right=2.5mm of t2] (k8) {};
    \node[bp key,right=1mm of k8] (t3) {编号假设};
    \node[bp math,anchor=west,minimum width=4.5mm,right=2.5mm of t3] (k9) {};
    \node[bp key,right=1mm of k9] (t4) {纯数学};
    \node[bp stub,anchor=west,minimum width=4.5mm,right=2.5mm of t4] (k10) {};
    \node[bp key,right=1mm of k10] (t5) {其他单元};
    \node[bp fold,anchor=west,minimum width=4.5mm,right=2.5mm of t5] (k11) {$+k$};
    \node[bp key,right=1mm of k11] {折叠};
    \draw[bp use] (0,-0.42) -- ++(6mm,0) coordinate (e1);
    \node[bp key] (u1) at ($(e1)+(1.2mm,0)$) {依赖（前提指向陈述）};
    \draw[bp asmuse] ($(u1.east)+(2.5mm,0)$) -- ++(6mm,0) coordinate (e2);
    \node[bp key] (u2) at ($(e2)+(1.2mm,0)$) {编号假设 $\to R_{\rm cond}$};
    \draw[bp neck] ($(u2.east)+(2.5mm,0)$) -- ++(6mm,0) coordinate (e3);
    \node[bp key] (u3) at ($(e3)+(1.2mm,0)$) {瓶颈依赖（主表）};
    \draw[bp use,bp long] ($(u3.east)+(2.5mm,0)$) -- ++(6mm,0) coordinate (e4);
    \node[bp key] (u4) at ($(e4)+(1.2mm,0)$) {跨两列以上};
    \draw[bp stubuse] ($(u4.east)+(2.5mm,0)$) -- ++(6mm,0) coordinate (e5);
    \node[bp key] at ($(e5)+(1.2mm,0)$) {来自其他单元};
  \end{scope}
\end{tikzpicture}}
\caption{第一部分（第~1--5~章）主要定理的依赖图（由 \texttt{\string\uses} 自动生成）：总表中的陈述与本部分的全部定理、它们在本部分内的直接前提及所用的编号假设（40 个陈述节点、4 个汇总节点）。框的颜色为严格性标签，虚线框为 $R_{\rm cond}$ 条件定理，橙色虚线框为编号假设，椭圆为猜想，双线框为 [V]，白色框为纯数学；灰色框代表本部分之外的单元（该单元中被直接使用的陈述合并为一个节点），虚线白框“$+k$”为按折叠规则省略的 $k$ 个次要陈述。边由前提指向陈述，是可达关系在所画节点上的传递约化（经由未画出陈述的依赖也画成一条边）；橙色虚线边为编号假设到 $R_{\rm cond}$ 陈述的依赖，粗红线为决定全阶论证地位的瓶颈依赖（取自主表“瓶颈”列）。}\label{fig:appP-dependency-b}
\end{figure}
